\clearpage
\thispagestyle{empty}
\refstepcounter{section}
\label{app:preliminary}
\begin{center}
\textbf{APPENDIX D}

\vspace{\baselineskip}
\textbf{Selected Preliminary Computational Evidence}
\end{center}
\vspace{3\baselineskip}

\subsection{Status, Purpose, and Provenance}

This appendix preserves selected completed exploratory findings that
motivated the proposed evaluation. These outputs were reported in the
pre-consultation full-study manuscript and have not been regenerated
under every newly specified proposal convention. They demonstrate
implementation feasibility and expose limitations; they are not the final
results of the proposed study.

The preliminary core used recovered reconstructed validation inputs and
a frozen formulation. The restored scenario inventory contained 136 runs.
The five controlled fixtures contain 6, 8, 8, 7, and 20 nodes, respectively.
The software topology and measured electrical loading are real-data
inputs, while their disturbance mappings or scenario injections remain
distinct from independent incident outcomes. A complete pre-consultation
PDF and the original chapter sources are retained in the accompanying
research archive.

\subsection{Controlled Behavior and Comparator Differences}

The baseline outlet states were 0.053760, 0.408000, 0.106814, 0.351372,
and 0.708640 for N1--N5. These values are comparative indices under the
declared fixtures, not failure probabilities.

\begin{table}[H]
\centering
\caption{Preliminary N5 comparator agreement with heterogeneous GBBRPM.}
\label{tab:prelim-comparators}
\small
\begin{tabularx}{\textwidth}{@{}Lrr@{}}
\toprule
Comparator & Spearman $\rho$ & Top-three Jaccard\\
\midrule
Local-only & 0.108627 & 0.0\\
Uniform $S=0.25$ & -0.042105 & 0.0\\
Uniform $S=0.50$ & 0.150376 & 0.0\\
Uniform $S=0.75$ & 0.696241 & 0.5\\
\bottomrule
\end{tabularx}
\end{table}

These differences show that the supplied relation conditions and
propagated context changed the ordering. Agreement with GBBRPM is not
a ground-truth accuracy measure, and the table does not establish that
GBBRPM is operationally superior to the comparison scores.

\subsection{Preliminary Property, Robustness, and Dependence Checks}

The prior report recorded 600 randomized property-test trials with no
reported failure in its boundedness, local-disturbance monotonicity,
$B=1$ boundary, and zero-susceptibility checks. This is finite
implementation evidence alongside the analytical arguments, not a proof
of all possible extensions or implementations.

\begin{table}[H]
\centering
\caption{Preliminary N5 robustness summary, 600 trials per perturbation level.}
\label{tab:prelim-robustness}
\small
\begin{tabularx}{\textwidth}{@{}l*{5}{Y}@{}}
\toprule
Perturbation & Mean Spearman & 5th percentile & 95th percentile & Mean top-three Jaccard & Mean top-20\% Jaccard\\
\midrule
$\pm5\%$ & 0.979190 & 0.962331 & 0.992481 & 0.744500 & 0.776000\\
$\pm10\%$ & 0.938267 & 0.886289 & 0.974436 & 0.546667 & 0.604873\\
$\pm20\%$ & 0.845552 & 0.697067 & 0.935338 & 0.388833 & 0.440823\\
\bottomrule
\end{tabularx}
\end{table}

These previously reported set conventions are kept distinguishable from
the explicit cutoff-tie protocol for planned runs. The top-20\% metric has
nominal $k=4$ for N5. Small perturbations produced broadly stable
ordering, while the highest-priority membership was more sensitive.

The ancestry-aware path-pruned diagnostic gave outlet-state differences
of 0.000000 for N3, 0.099372 for N4, and 0.322637 for N5 relative to
ordinary aggregation. This is sensitivity to a declared counterfactual,
not a validated estimate of true physical bias. It demonstrates why a
bounded output and stable ranking do not establish dependence correction.

\subsection{Preliminary Runtime Evidence}

The existing sparse-DAG experiment reported median timings from
1.628378 ms for 20 nodes and 60 edges to 16.650025 ms for 1000 nodes
and 3000 edges, after five warm-ups and 30 repetitions.
These implementation-specific timings do not provide a matched speed
comparison with a hydraulic simulator, memory evidence, or a breakdown
of validation, propagation, and sorting. The planned cost protocol will
supply those distinctions rather than infer them from this table range.

\subsection{Preliminary Hydraulic-Reference Disagreement}

The prior 12-scenario SWMM comparison reported the following:
\begin{table}[H]
\centering
\caption{Preliminary comparison with simulated hydraulic worsening.}
\label{tab:prelim-swmm}
\small
\begin{tabularx}{\textwidth}{@{}L*{3}{Y}@{}}
\toprule
Hydraulic outcome & Mean Spearman & Top-20\% Jaccard & Positive worsening outside downstream scope\\
\midrule
Maximum depth & -0.416 & 0.119 & 71.2\%\\
Flooding volume & -0.485 & 0.000 & 40.5\%\\
\bottomrule
\end{tabularx}
\end{table}

The negative associations and low priority-set overlaps do not support
hydraulic-outcome ranking equivalence under this protocol. Outside-scope
worsening is consistent with mechanisms omitted by the downstream-only
DAG index, but the comparison does not isolate a causal hydraulic
mechanism in each scenario. Neither disagreement nor low computation
cost validates drainage inspection or maintenance utility.

\subsection{Preliminary Software and Electrical Instantiation}

The Express 4.18.2 case used 71 packages, 128 directed relations, and
seven advisory-backed local disturbances. Root risk increased from
0.5000 at $\tau=0$ to 0.9975 at $\tau=1$, and its rank became first
by $\tau=0.50$. Appendix~\ref{app:software-worked-example} traces the
intermediate setting. Shared \texttt{qs} and \texttt{send} paths may
inflate the root state; their contribution is not quantified by the
stable ranking.

The observable Component 2 case used 12 nodes, 11 edges, and loading
derived from 24 aligned phasor snapshots. The 43 controlled scenarios
tested severity, location, load scaling, transmission, and multiple local
sources. The root-disturbance case $B_{\mathrm{C2\_N01}}=0.75$ gave
risk sum approximately 1.186. Appendix~\ref{app:electrical-worked-example}
provides its trace. These are injected-disturbance experiments on
normal-operation measurements, not independent event-response prediction.

\subsection{Preliminary Prototype Feasibility}

The prior implementation report recorded a controlled-fixture inventory,
frozen N5 reproduction, valid imported-graph evaluation, invalid-input
rejection, and a production frontend build. Those checks support
implementation feasibility within the tested configuration. The planned
study will rerun and extend engine/prototype consistency checks; it will
not infer usability, security, field applicability, or operational
readiness from the preliminary functional tests.
